\section{Discussion and limitations}
\label{sec:discussion}

\paragraph{Practical guidance, with its scope}
For a practitioner with a small labeled set, a large unlabeled set and a contrastive pretraining pipeline,
our results support three rules.
\begin{enumerate}[leftmargin=1.6em]
\item \emph{Given enough pretraining, do not spend scarce labels in a \supcon-style loss from the start; keep
them for fine-tuning or for semi-supervised training.} If the encoder will be used frozen with a linear head, a brief \supcon stage
after label-free pretraining is a good use of them; it does not help when the whole network is fine-tuned
or trained with FixMatch afterwards. If semi-supervised training follows pretraining, the choice of
pretraining matters little, and with a fixed compute budget semi-supervised training alone was as good as
any pretrained start we tried.
\item \emph{If pretraining compute is very limited, labels in the loss are the faster route, unless they are very
few.} Below roughly 5,000 steps, \supcon on the labeled subset is ahead under every protocol we ran with 5--50\%
of the CIFAR-10 labels; with 1\% (50 images per class) it is already behind at 2,425 steps
(\cref{sec:fractions}). At 9,700 steps the answer depends on how the labels are split: with all 5,000 labels in
the loss, fine-tuning favours \simclr; with 1,000 held out to choose the recipe, \supcon is
\R{cmp:simclr-holdsupcon:ftt:e100:diff}\pt ahead and \union level (\cref{sec:protocol}). By 19,400 steps
\simclr leads either way.
\item \emph{Evaluate with the labels you have.} Report a label-budget protocol next to the all-label probe,
say how its hyper-parameters were chosen (an encoder that used the labels cannot be tuned on them; hold some
out before pretraining), and state budgets in steps or processed images.
\end{enumerate}
These rules are supported for class-balanced data, unlabeled images from the same or a broader distribution
(CIFAR, STL-10, Tiny-ImageNet), ResNet-18 and ResNet-50 encoders, label fractions around 1--20\% and budgets up to
\R{long:steps} steps. With all labels \supcon is the better pretraining method \citep{khosla2020supcon};
in our grid the label fraction from which it overtakes \simclr grows with the budget, to
\R{grid:supcon:lpall:e400:be} under the probes at 38,800 steps, and after fine-tuning \simclr is still ahead at
50\% of the labels at that budget (\cref{sec:fractions}).

\paragraph{Limitations}
\emph{Scale.} All experiments use images of at most 96 px, at most 105,000 training images and ResNets. We
did not run ImageNet-scale data or vision transformers, and the budgets at which the crossover occurs will
differ there; for the ResNet-50 it already comes later than for the ResNet-18.

\emph{Balanced, curated labels.} The labeled subsets are class-balanced. Self-supervised pretraining is
reported to be more robust to class imbalance than supervised pretraining
\citep{liu2022self,yang2020rethinking}, which would favour our conclusion, but we did not test it.

\emph{Unlabeled data.} On CIFAR and Tiny-ImageNet the unlabeled images come from the same distribution as the
labeled ones.
STL-10 relaxes this, but its unlabeled set is still close to the labeled classes. With unlabeled data from
a different domain, the value of the ``all images'' column can be lower \citep{cole2022when}.

\emph{The label-using arms.} \union is the simplest way of mixing labeled and unlabeled images in a
contrastive loss. Methods that add pseudo-labels or a class-conditional instance term
\citep{li2021comatch,zheng2022simmatch,chen2022perfectly} address exactly the weakness we measure, and our
negative result for \union does not extend to them.

\emph{Fine-tuning.} The recipe search of \cref{sec:protocol} covers 72 recipes for one architecture, 120 for
the label-using encoders at 19,400 steps, where widening the grid changed the best accuracy by at most a few
tenths of a point. Augmentation policies such as RandAugment \citep{cubuk2020randaugment} could still move the
fine-tuning numbers of all encoders.

\emph{Statistics.} Three to five seeds give a coarse estimate of the between-seed variance; several
secondary results rest on one seed and are marked. The coverage analysis is an intervention on the pool,
but its link to the plateau of \supcon remains correlational.

\section{Conclusion}
\label{sec:conclusion}
We asked whether scarce labels should go inside a contrastive loss or be used after label-free pretraining,
and removed the two confounds that make the question hard: unequal compute and evaluation with labels that
are not available. At equal steps, labels inside the loss give a fast start and an early plateau, and
label-free pretraining on all images overtakes them. Which method looks better at the crossover depends on
the protocol: the all-label probe favours the label-using methods, and the protocols that respect the label
budget favour label-free pretraining when all labels went into pretraining, strongly under supervised fine-tuning
with a recipe a practitioner can choose and weakly under semi-supervised fine-tuning. Holding out a fifth of the
labels to choose the fine-tuning recipe changes this at the crossover: \supcon is then ahead and \union level. With longer training label-free pretraining
leads under every probe and fine-tuning recipe, while after semi-supervised fine-tuning the encoders are
equivalent. The geometry of the features is consistent with this: the labels tighten the labeled images
far more than any other images. Within contrastive pretraining the labels
are best spent last; across all methods we tested, they are best spent in semi-supervised training.
